\scanpage{33}
% archive foliation: 32
En fait, on trouve un diagramme
can. commutatif
\diagram{\begin{tikzcd}[column sep=small,row sep=huge]
&&A\arrow[r,phantom,"\subset"]\arrow[dll,"\operatorname{id}"']\arrow[drr,"\operatorname{id}"]&
C\arrow[dll,"\rho","\sim"']\arrow[drr,"\rho'","\sim"']\\
A\arrow[r,phantom,"\subset"]&\widehat A_{\unsure{\mathcal{T}|\tau}}\arrow[r,phantom,"\subset"]&\widehat A&&
A\arrow[r]& (\widehat{A^\circ}_{\unsure{\mathcal{T}|\tau}'})^\circ\arrow[r,phantom,"\subset"]&
(\widehat{A^\circ})^\circ
\end{tikzcd}}
(resp. sur $\widehat A_{ab}$ et $\widehat{A^\circ}_{ab}$)
On trouve en particulier une
équivalence remarquable
\[
\widehat A_{\unsure{\mathcal{T}|\tau}}\simeq(\widehat{A^\circ}_{\unsure{\mathcal{T}|\tau}'})^\circ
\]
qui peut se définir intrinsèque-
ment \inserted{en termes de $A,\unsure{\mathcal{T}|\tau},\unsure{\mathcal{T}|\tau}'$}, chaque fois qu'on a une
sous-cat. $D$ de $\widehat A$ (resp. une sous-cat.
additive de $\widehat A_{ab}$), stable \inserted{dans $\widehat A$} par
$\varprojlim$ de type $\unsure{\mathcal{T}|\tau}$, et \inserted{où existent les} $\varinjlim$ de type $\unsure{\mathcal{T}|\mathcal{T}'}$
(pas néc. les mêmes que dans $\widehat A$!)
et dont chaque objet est à la fois
$\varprojlim$ \inserted{de type $\unsure{\mathcal{T}|\tau}$} dans $\unsure{C|D}$ de foncteurs repr.,
et $\varinjlim$ \inserted{de type $\unsure{\mathcal{T}|\tau}'$} dans $\unsure{C|D}$ de tels foncteurs.

Si $A$ est réduite à un seul objet \inserted{$A$} et
dans le cas resp. ($C$ additive)
$\unsure{\mathcal{T}|\tau}=\unsure{\mathcal{T}|\tau}'=$ \deleted{\unsure{ens.}} catégories discrètes finies,
on trouve une antiéquivalence
entre modules libres de t.f. à g. et
